\section{Método de selección de literatura}

La revisión se construyó a partir de un corpus curado de trece artículos científicos disponibles localmente en la carpeta \path{context/papers/}. El corpus incluye artículos de revista y conferencias publicadas entre 2016 y 2025, y fue delimitado para cubrir cuatro focos complementarios: detección general de trastornos mentales con PLN, detección de ideación suicida en redes sociales, enfoques de depresión fundamentados clínicamente en DSM-5 y desarrollos recientes basados en aprendizaje profundo, modelos basados en transformers y LLMs.

Los criterios de inclusión fueron los siguientes: (i) abordar explícitamente salud mental, depresión o ideación suicida en texto de redes sociales; (ii) utilizar o discutir metodologías de PLN, aprendizaje automático, aprendizaje profundo, modelos basados en transformers o LLMs; (iii) tener relevancia directa para Reddit o para tareas comparables de detección en redes sociales; y (iv) aportar evidencia para la dimensión de análisis seleccionada, es decir, efectividad predictiva y/o interpretabilidad clínica. Bajo estos criterios, se privilegiaron artículos que describen el origen de los datos, la representación textual, el modelo computacional, las métricas reportadas y las implicaciones prácticas de sus resultados.

La literatura seleccionada se analizó de forma comparativa y no como una lista de resúmenes independientes. Para ello, cada artículo fue caracterizado mediante los siguientes atributos: problema de salud mental abordado, fuente de datos, representación textual, modelo, métricas reportadas, contribución principal, limitaciones y aportación específica a la efectividad predictiva o a la interpretabilidad clínica. Este esquema permite alinear la revisión con la rúbrica de evaluación del curso, porque fuerza a explicitar no solo qué hace cada trabajo, sino también por qué su contribución resulta pertinente para el reto.

Finalmente, los trabajos fueron agrupados por enfoque de investigación. Esta organización facilita distinguir entre estudios panorámicos, modelos tempranos basados en rasgos lingüísticos, enfoques clínicamente anclados en DSM-5 y modelos contemporáneos que priorizan potencia predictiva mediante arquitecturas profundas o LLMs. La Tabla~\ref{tab:state-of-the-art} sintetiza el corpus completo y funciona como referencia transversal para las secciones analíticas posteriores.
